% GENERATED FILE. Edit canonical lessons, then run npm run generate:books.
% canonical-source-hash: fnv1a64:5ec977fec28d742e
% canonical-lessons: SA-C50-sound, SA-C50-question, SA-C50-voice, SA-C50-song, SA-C50-verse

\chapter{Sound and Song}
\label{ch:sa-sound-and-song}
\hlchaptermodality{50}

\begin{chapteropening}
\textbf{By the end of this chapter:} \emph{I can name a sound, a question, the voice, a song and a verse in Sanskrit.}

Complete SA-C50-verse and demonstrate the chapter promise.
\end{chapteropening}

\section[śabdaḥ]{\sk{शब्दः} (śabdaḥ) — a sound}

\label{lesson:SA-C50-sound}



\begin{quote}
Before the new word: what did \emph{aho} mean?
\end{quote}

\subsection*{You'll want to know: \sk{शब्दः}}
\textbf{\sk{शब्दः}} (\emph{śabdaḥ}) — \textquotedblleft{}a sound\textquotedblright{}.

\textbf{\sk{शब्दः}} is a sound, and by an easy stretch a word said aloud.

This one has been handed to you three times without ever being named. \emph{kāla-śabdāḥ} were the time ones; \emph{ka-śabdāḥ} were the question ones; \emph{saṅkhyā-śabdāḥ} were the number ones. The \emph{-śabdāḥ} ending each of those three names is this very word, said of more than one thing at a time.

So you have been reading it since long before this page. Now the piece itself is yours. Ends in \textbf{-\sk{अः}}.

\begin{grammarlens}[title={What you've built}]
The piece three earlier names were built out of.
\end{grammarlens}

\subsection*{Your turn}
Say these aloud:

\begin{itemize}
\raggedright
  \item \emph{śabdaḥ}
  \item it once more, slowly
  \item \emph{aho}, then \emph{śabdaḥ}
\end{itemize}

\begin{tcolorbox}[breakable,colback=teal!4,colframe=teal!35!black,title={Before you move on}]
What does \sk{शब्दः} mean? (\textquotedblleft{}a sound\textquotedblright{}.) Say it once more.
\end{tcolorbox}
\section[praśnaḥ]{\sk{प्रश्नः} (praśnaḥ) — a question}

\label{lesson:SA-C50-question}



\begin{quote}
Before the new word: what did \emph{śabdaḥ} mean?
\end{quote}

\subsection*{You'll want to know: \sk{प्रश्नः}}
\textbf{\sk{प्रश्नः}} (\emph{praśnaḥ}) — \textquotedblleft{}a question\textquotedblright{}.

\textbf{\sk{प्रश्नः}} is a question: the thing asked, rather than the asking of it.

You already say \sk{पृच्छति}, \textquotedblleft{}asks\textquotedblright{}. \textbf{\sk{प्रश्नः}} stands on that same root, held still as a thing instead of running as an action — and the pair is worth saying together, because Sanskrit does this everywhere.

The root reaches a long way west: Latin \emph{precārī}, \textquotedblleft{}to beg\textquotedblright{}, behind English \emph{pray} and \emph{precarious}, and German \emph{fragen}, \textquotedblleft{}to ask\textquotedblright{}. Ends in \textbf{-\sk{अः}}.

\begin{grammarlens}[title={What you've built}]
The asking, and the thing asked, from one root.
\end{grammarlens}

\subsection*{Your turn}
Say these aloud:

\begin{itemize}
\raggedright
  \item \emph{praśnaḥ}
  \item it once more, slowly
  \item \emph{śabdaḥ}, then \emph{praśnaḥ}
\end{itemize}

\begin{tcolorbox}[breakable,colback=teal!4,colframe=teal!35!black,title={Before you move on}]
What does \sk{प्रश्नः} mean? (\textquotedblleft{}a question\textquotedblright{}.) Say it once more.
\end{tcolorbox}
\section[vāṇī]{\sk{वाणी} (vāṇī) — the voice}

\label{lesson:SA-C50-voice}



\begin{quote}
Before the new word: what did \emph{praśnaḥ} mean?
\end{quote}

\subsection*{You'll want to know: \sk{वाणी}}
\textbf{\sk{वाणी}} (\emph{vāṇī}) — \textquotedblleft{}the voice\textquotedblright{}.

\textbf{\sk{वाणी}} is the voice: speech as a thing heard, rather than speech as a thing meant.

It rests on a piece meaning \textquotedblleft{}to sound\textquotedblright{}. Beside \sk{शब्दः}, which is any sound at all, \textbf{\sk{वाणी}} is the sound a person makes on purpose.

Capitalised, \emph{Vāṇī} is a name of the goddess of speech, so the word carries a good deal more weight in poetry than its plain sense suggests. Feminine, ending in \emph{-ī}.

\begin{grammarlens}[title={What you've built}]
A sound narrowed down to the human one.
\end{grammarlens}

\subsection*{Your turn}
Say these aloud:

\begin{itemize}
\raggedright
  \item \emph{vāṇī}
  \item it once more, slowly
  \item \emph{praśnaḥ}, then \emph{vāṇī}
\end{itemize}

\begin{tcolorbox}[breakable,colback=teal!4,colframe=teal!35!black,title={Before you move on}]
What does \sk{वाणी} mean? (\textquotedblleft{}the voice\textquotedblright{}.) Say it once more.
\end{tcolorbox}
\section[gītam]{\sk{गीतम्} (gītam) — a song}

\label{lesson:SA-C50-song}



\begin{quote}
Before the new word: what did \emph{vāṇī} mean?
\end{quote}

\subsection*{You'll want to know: \sk{गीतम्}}
\textbf{\sk{गीतम्}} (\emph{gītam}) — \textquotedblleft{}a song\textquotedblright{}.

\textbf{\sk{गीतम्}} is a song: the thing sung.

It stands on \emph{gai}, \textquotedblleft{}to sing\textquotedblright{}. The \emph{Bhagavad Gītā} is this same word wearing a feminine ending — \textquotedblleft{}the song\textquotedblright{}, and nothing more mysterious than that.

No secure western cousin has been agreed for it, but it kept its ground at home: Hindi \emph{gīt} and Bengali \emph{gīt} are still the everyday word for a song. Neuter, ending in \textbf{-\sk{अम्}}.

\begin{grammarlens}[title={What you've built}]
What the voice does when it is not talking.
\end{grammarlens}

\subsection*{Your turn}
Say these aloud:

\begin{itemize}
\raggedright
  \item \emph{gītam}
  \item it once more, slowly
  \item \emph{vāṇī}, then \emph{gītam}
\end{itemize}

\begin{tcolorbox}[breakable,colback=teal!4,colframe=teal!35!black,title={Before you move on}]
What does \sk{गीतम्} mean? (\textquotedblleft{}a song\textquotedblright{}.) Say it once more.
\end{tcolorbox}
\section[ślokaḥ]{\sk{श्लोकः} (ślokaḥ) — a verse}

\label{lesson:SA-C50-verse}



\begin{quote}
Before the last word of the chapter: what did \emph{vāṇī} mean, and what did \emph{gītam} mean?
\end{quote}

\subsection*{You'll want to know: \sk{श्लोकः}}
\textbf{\sk{श्लोकः}} (\emph{ślokaḥ}) — \textquotedblleft{}a verse\textquotedblright{}.

\textbf{\sk{श्लोकः}} is a verse: two lines, balanced against each other.

Its older sense is \textquotedblleft{}sound, and what is heard of a person\textquotedblright{} — fame, in other words. So a verse is named as a thing that carries a sound outward. Beside \sk{शब्दः}, which is any sound, \textbf{\sk{श्लोकः}} is a sound cut to a shape that can be remembered.

Nearly the whole of the long Sanskrit story-poems is in this one shape, which is how the word came to mean the shape as well as the fame. Ends in \textbf{-\sk{अः}}.

\begin{grammarlens}[title={What you've built}]
Five things the ear takes in: a sound, a question, the voice, a song, a verse.
\end{grammarlens}

\subsection*{Your turn}
Say these aloud:

\begin{itemize}
\raggedright
  \item \emph{ślokaḥ}
  \item it once more, slowly
  \item all five in order, then \emph{śabdaḥ} and \emph{ślokaḥ} together
\end{itemize}

\begin{tcolorbox}[breakable,colback=teal!4,colframe=teal!35!black,title={Before you move on}]
What does \sk{श्लोकः} mean? (\textquotedblleft{}a verse\textquotedblright{}.) Say it once more.
\end{tcolorbox}
